\chapter{Einleitung}
\label{kap:einleitung}

Die Einleitung führt in das Thema ein: Relevanz, Problemstellung, Zielsetzung,
Methode und Aufbau der Arbeit. Zwischen zwei Überschriften steht immer etwas Text.

\section{Zitieren nach Harvard}

Zitiert wird immer mit \verb|\autocite|. Ein direktes Zitat mit Seitenangabe:
\enquote{Wörtlich übernommener Text} \autocite[12]{mustermann2023}. Eine
sinngemäße Übernahme bekommt ein \enquote{vgl.} \autocite[vgl.][45--47]{schmidt2021}.
Steht der Name im Satz, passt \verb|\textcite|: \textcite{weber2024} zeigt, dass
\dots{} Internetquellen werden genauso zitiert \autocite{beispielblog2025}.

\section{Abkürzungen und Verweise}

Beim ersten Vorkommen wird eine Abkürzung ausgeschrieben: \ac{edi}. Danach steht nur
noch \ac{edi}. Auf Kapitel, Abbildungen und Tabellen wird automatisch verwiesen,
etwa auf \autoref{kap:grundlagen}, \autoref{abb:beispiel} oder \autoref{tab:beispiel}.
Zwischen Abkürzungsteilen steht ein schmales geschütztes Leerzeichen: z.\,B., d.\,h.,
u.\,a.; vor Einheiten ein geschütztes Leerzeichen: 100~mm.
